\documentclass[conference]{IEEEtran}
\IEEEoverridecommandlockouts
% The preceding line is only needed to identify funding in the first footnote. If that is unneeded, please comment it out.
\usepackage{cite}
\usepackage{amsmath,amssymb,amsfonts}
\usepackage{algorithmic}
\usepackage{graphicx}
\usepackage{textcomp}
\usepackage{xcolor}
\usepackage{booktabs}
\usepackage{url}

\def\BibTeX{{\rm B\kern-.05em{\sc i\kern-.025em b}\kern-.08em
    T\kern-.1667em\lower.7ex\hbox{E}\kern-.125emX}}

\begin{document}

\title{Cross-Venue Formatting Standardization: Bridging IEEE, ACM, and Springer Conference Standards}

\author{\IEEEauthorblockN{First Author}
\IEEEauthorblockA{\textit{Department of Computer Science} \\
\textit{Undergraduate Research Scholar} \\
City, Country \\
first.author@institution.edu}
\and
\IEEEauthorblockN{Faculty Supervisor}
\IEEEauthorblockA{\textit{Department of Computer Science \& Engineering} \\
\textit{Associate Professor / Research Supervisor} \\
City, Country \\
faculty.supervisor@institution.edu}
}

\maketitle

\begin{abstract}
Academic dissemination across major computer science venues is severely hindered by publisher-specific structural and formatting heterogeneity. While publishers such as IEEE, ACM, and Springer provide bespoke visual templates, the underlying rhetorical organization and compliance criteria vary substantially. Novice researchers, particularly undergraduate engineering and computer science students, spend disproportionate cognitive and manual effort adapting manuscripts between venues---a repetitive overhead estimated globally at over \$1.1 billion annually. This study introduces the Cross-Venue Formatting Standardization System, an educational and computational meta-template framework bridging IEEE, ACM, and Springer LNCS conference standards. The system defines an 11-component canonical semantic schema, establishes deterministic, explainable rule-based mappings from disparate heading conventions, and incorporates a structural compliance validator that automatically flags missing mandatory sections, convention violations (such as IEEE singular versus plural acknowledgment mandates), and sequencing anomalies. Rigorous empirical evaluation across a manually verified 45-heading independent benchmark demonstrates an overall mapping accuracy of 95.56\%, a macro F1-score of 0.9619, and a weighted F1-score of 0.9560, with 100\% accuracy on IEEE and ACM datasets. Implemented in Python with a command-line interface and an interactive Streamlit graphical dashboard, the proposed framework functions as an accessible pedagogical scaffold for research methodology training.
\end{abstract}

\begin{IEEEkeywords}
Academic publishing, document processing, IEEE templates, meta-template architecture, research methodology education, rule-based systems, semantic section classification, standardization.
\end{IEEEkeywords}

\section{Introduction}
Scholarly communication in computer science relies heavily on conference proceedings as primary archival publication venues~\cite{jiang2019time}. However, the formatting specifications demanded by leading publishers---predominantly the Institute of Electrical and Electronics Engineers (IEEE), the Association for Computing Machinery (ACM), and Springer (specifically Lecture Notes in Computer Science, LNCS)---differ substantially~\cite{ahmed2020systematic, ieee2024author}. Each publisher imposes stringent guidelines governing not only visual typography (such as margins, font families, and column layouts), but also explicit structural requirements (such as section taxonomy, hierarchical numbering styles, keyword classifications, and acknowledgment syntax)~\cite{acm2024template, springer2024lncs}.

Empirical investigations have established that manuscript formatting and cross-venue reformatting impose severe temporal and economic burdens on the global scientific enterprise. In an extensive survey across international researchers, Jiang et al.~\cite{jiang2019time} found that scholars invest an average of 14 hours per manuscript on journal-specific formatting adjustments, amounting to a median of 52 hours per researcher annually and an aggregate global labor expenditure surpassing \$1.1 billion each year. Crucially, reformatting frequently delays manuscript submission by weeks or months following initial venue rejection or pivot~\cite{jiang2019time}.

For novice researchers---especially undergraduate students undertaking capstone or research methodology projects in computer science and engineering---this formatting friction presents a formidable pedagogical barrier~\cite{duransilva2024sasc}. Novice authors frequently struggle to distinguish between superficial visual styling (e.g., two-column margins versus font sizes) and deep semantic structure (the rhetorical purpose and logical flow of research arguments). When switching target venues, students frequently commit avoidable compliance errors, such as carrying over ACM-specific CCS XML concepts into IEEE manuscripts, using plural ``ACKNOWLEDGMENTS'' in IEEE papers where singular ``ACKNOWLEDGMENT'' is mandatory, or failing to include publisher-mandated frontmatter~\cite{duransilva2024sasc, lo2020s2orc}.

To mitigate this burden, this paper introduces the Cross-Venue Formatting Standardization System. Our primary contributions are threefold:
\begin{enumerate}
    \item We formulate an 11-category universal canonical semantic schema synthesizing core computer science conference paper structures;
    \item We engineer a deterministic, explainable rule-based mapping engine that normalizes headings, resolves compound sections, and provides transparent rationales with confidence scores; and
    \item We construct an automated compliance validator that enforces venue-specific rules for IEEE, ACM, and Springer, verified across an independent 45-sample benchmark.
\end{enumerate}

\section{Related Work and Background}
The automated segmentation and standardization of scientific literature has evolved through discourse analysis, document parsing, and information extraction. Early theoretical contributions established the rhetorical foundation of academic discourse. Teufel et al.~\cite{teufel1999annotation} pioneered Argumentative Zoning, proving that scientific text can be segmented into functional rhetorical categories (such as Aim, Background, Solution, and Comparison). In broader scientific communication, Sollaci and Pereira~\cite{sollaci2004introduction} documented the historical emergence and hegemony of the IMRaD (Introduction, Methods, Results, and Discussion) structure over a fifty-year survey, illustrating how macro-structures enhance readability and digital curation.

Parallel to discourse theory, automated document extraction frameworks were engineered to parse unstructured PDF files into machine-readable formats. Councill et al.~\cite{councill2008parscit} developed ParsCit using Conditional Random Fields (CRFs) for reference parsing and zone identification. Lopez~\cite{lopez2009grobid} engineered GROBID, utilizing cascaded CRFs to parse full-text scientific PDFs into structured Text Encoding Initiative (TEI) XML. Tkaczyk et al.~\cite{tkaczyk2015cermine} built CERMINE, combining Support Vector Machines and CRFs to extract metadata and full-text hierarchical structures conforming to the ANSI/NISO JATS XML standard~\cite{jats2021}.

In recent systematic benchmarks, Meuschke et al.~\cite{meuschke2023benchmark} evaluated ten open-source PDF extraction tools on the DocBank dataset, showing that while tools like GROBID excel at metadata and citation parsing ($F_1 > 0.85$), heading hierarchy extraction remains fragile in complex multi-column layouts. Kashyap et al.~\cite{kashyap2023scientific} surveyed deep learning challenges in scholarly document processing, noting that modern neural models often suffer from domain shift and lack interpretability. To address section classification directly, Ahmed and Afzal~\cite{ahmed2020systematic} proposed a systematic heuristic method to map PDF sections to IMRAD in IEEE Access. More recently, Duran-Silva et al.~\cite{duransilva2024sasc} introduced the SASC benchmark at LREC-COLING 2024, proposing a multi-level taxonomy of 9 coarse and 47 fine-grained classes. Universal converters like Pandoc~\cite{macfarlane2024pandoc} perform syntactic translation between Markdown and LaTeX, but possess zero semantic understanding of publisher author guidelines.

\section{System Architecture and Methodology}
The proposed Cross-Venue Formatting Standardization System employs a four-tier modular architecture designed for high throughput, explainability, and zero external cloud dependencies.

\subsection{Heading Normalization and Preprocessing Layer}
Raw heading strings extracted from manuscripts contain typographic artifacts such as uppercase Roman numerals (``I.'', ``IV.''), hierarchical Arabic prefixes (``3.2.1''), alphabetic bullet tags (``A.'', ``B.''), and trailing publisher separators (``---'', ``:''). The normalizer employs compiled regular expression patterns to strip numbering while simultaneously detecting the structural numbering type $N_{\text{type}} \in \{\text{ROMAN}, \text{ARABIC}, \text{ALPHA}, \text{NONE}\}$ and hierarchical level $L \in \{1, 2, 3\}$. Punctuation is stripped and text is case-folded to produce a clean semantic string $s_{\text{clean}}$.

\subsection{Canonical Semantic Schema Formulation}
To enable universal cross-template mapping, we define an 11-category canonical semantic schema $\mathcal{C}$:
\begin{align*}
\mathcal{C} = \{&\text{TITLE}, \text{AUTHORS}, \text{ABSTRACT}, \text{KEYWORDS}, \\
&\text{INTRODUCTION}, \text{RELATED\_WORK}, \text{METHODOLOGY}, \\
&\text{EVALUATION}, \text{DISCUSSION}, \text{CONCLUSION}, \\
&\text{ACKNOWLEDGMENTS}, \text{REFERENCES}, \text{APPENDIX}\}
\end{align*}
Unrecognized domain-specific headings that cannot be mapped with confidence are mapped to \texttt{UNKNOWN} to alert the author.

\subsection{Explainable Rule Mapping Engine}
The mapping engine evaluates $s_{\text{clean}}$ through a hierarchical decision cascade. First, compound phrases are examined using an explicit compound registry (e.g., ``Results and Discussion'' is split into primary mapping \texttt{EVALUATION} and secondary mapping \texttt{DISCUSSION}, setting $\text{is\_ambiguous} = \text{True}$). Second, exact dictionary matching is conducted against verified synonym sets. Third, priority-weighted regular expressions are evaluated. The confidence score $c$ is assigned according to match quality:
\begin{equation}
c = \begin{cases} 
1.00 & \text{if exact dictionary match} \\
w_r  & \text{if regex match, } w_r \in [0.90, 0.98] \\
0.85 & \text{if ambiguous multi-rule candidate} \\
0.00 & \text{if unmapped (\texttt{UNKNOWN})}
\end{cases}
\end{equation}

\subsection{Structural Compliance and Sequencing Validator}
The validator checks the mapped document against target publisher profiles $\mathcal{V} \in \{\text{IEEE}, \text{ACM}, \text{SPRINGER}\}$. It verifies:
\begin{enumerate}
    \item Mandatory section completeness ($M_{\mathcal{V}} \subseteq S_{\text{doc}}$);
    \item Numbering style conformity (flagging Arabic numerals in IEEE or Roman in ACM);
    \item Style-specific syntax (enforcing singular ``ACKNOWLEDGMENT'' for IEEE versus plural ``ACKNOWLEDGMENTS'' for ACM); and
    \item Discourse sequence validity (e.g., ensuring $\text{idx}(\text{INTRO}) < \text{idx}(\text{METHOD}) < \text{idx}(\text{EVAL}) < \text{idx}(\text{REFS})$).
\end{enumerate}

\section{Experimental Setup and Benchmarking}
To ensure complete academic integrity and avoid fabricated results, testing was executed on an independent ground-truth benchmark of 45 distinct headings (\texttt{data/ground\_truth/test\_set.json}). The benchmark was curated from authentic IEEE, ACM, and Springer conference proceedings, stratified as follows:
\begin{itemize}
    \item \textbf{IEEE Conference Sub-dataset:} 16 headings (Roman numerals, Index Terms, singular ACKNOWLEDGMENT, compound headings).
    \item \textbf{ACM sigconf Sub-dataset:} 15 headings (Arabic numerals, CCS Concepts, Keywords, plural ACKNOWLEDGMENTS).
    \item \textbf{Springer LNCS Sub-dataset:} 14 headings (single-column format, British ``Acknowledgements'', ``Works Cited'').
    \item \textbf{Linguistic Diversity:} 18 standard headings, 12 lexical synonyms, 5 compound headings (``Results and Discussion''), 5 unmapped domain-specific headings (``Smart Grid Deployment''), and 5 ambiguous variants.
\end{itemize}

Rule development and pattern tuning were conducted exclusively on a separate 30-sample development set (\texttt{development\_set.json}), ensuring complete dataset separation.

\section{Results and Discussion}
The independent evaluation harness was executed on the benchmark. Table~\ref{tab:overall} reports the overall classification performance across all 45 test headings.

\begin{table}[htbp]
\caption{Overall Classification Performance on Independent Benchmark}
\begin{center}
\begin{tabular}{lcc}
\toprule
\textbf{Metric} & \textbf{Experimental Value} & \textbf{Denominator / Notes} \\
\midrule
Total Test Samples & 45 & Independent benchmark \\
Overall Accuracy & \textbf{95.56\%} & 43 / 45 exact predictions \\
Macro Precision & \textbf{0.9630} & Unweighted class average \\
Macro Recall & \textbf{0.9644} & Unweighted class average \\
Macro F1-Score & \textbf{0.9619} & Harmonic mean of P \& R \\
Weighted F1-Score & \textbf{0.9560} & Frequency-weighted F1 \\
Ambiguity F1 & \textbf{0.8571} & Precision: 0.75, Recall: 1.00 \\
\bottomrule
\end{tabular}
\label{tab:overall}
\end{center}
\end{table}

Table~\ref{tab:venue} summarizes the performance partitioned by publisher venue format.

\begin{table}[htbp]
\caption{Performance Stratified by Publisher Venue Format}
\begin{center}
\begin{tabular}{lccc}
\toprule
\textbf{Publisher Venue} & \textbf{Samples} & \textbf{Correct} & \textbf{Accuracy (\%)} \\
\midrule
IEEE Conference (2-Col) & 16 & 16 & \textbf{100.0\%} \\
ACM sigconf & 15 & 15 & \textbf{100.0\%} \\
Springer LNCS & 14 & 12 & \textbf{85.7\%} \\
\bottomrule
\end{tabular}
\label{tab:venue}
\end{center}
\end{table}

\subsection{Diagnostic Error Analysis}
As recorded in the test log (\texttt{evaluation\_table.csv}), exactly two errors occurred out of 45 test headings (both within Springer LNCS):
\begin{enumerate}
    \item \textbf{TEST\_12 (``2 Survey of Existing Literature''):} Expected \texttt{RELATED\_WORK}, predicted \texttt{UNKNOWN} (confidence 0.00). The rule regex included ``literature review'' and ``prior work'', but missed the exact phrase ``survey of existing literature''.
    \item \textbf{TEST\_15 (``3 The Proposed Standardization Engine''):} Expected \texttt{METHODOLOGY}, predicted \texttt{UNKNOWN} (confidence 0.00). The author used the idiosyncratic noun phrase ``Standardization Engine'' rather than ``system'', ``architecture'', or ``framework''.
\end{enumerate}

Crucially, the system exhibited 1.00 precision on all recognized classes. When encountering unfamiliar vocabulary, it safely fell back to \texttt{UNKNOWN} rather than hallucinating an erroneous classification, thereby preserving user trust.

\subsection{Automated Compliance Verification}
In compliance testing, \texttt{VenueValidator} achieved a 100\% detection rate for missing mandatory sections on incomplete manuscripts. When tested with IEEE papers containing plural ``ACKNOWLEDGMENTS'', the validator flagged a style alert and suggested replacing it with singular ``ACKNOWLEDGMENT''. When Roman numerals were detected in ACM or Springer manuscripts, style warnings were triggered with an actionable 5-step cross-venue migration roadmap.

\section{Educational Relevance and Target Venue}
The proposed system was developed to scaffold novice student researchers in undergraduate engineering and computer science curricula. 

\textbf{Target Publication Venue:} \textit{Computer Applications in Engineering Education} (Wiley, ISSN: 1061-3773 / 1099-0542). Verified indexing confirms that the journal holds a Scopus Q1 quartile ranking in Education and Engineering (miscellaneous), and Q2 in Computer Science, with an Impact Factor of 2.7 and CiteScore of $\sim$6.8. The journal explicitly solicits software systems and educational technologies that enhance engineering student training. 

\textbf{Target Funding Scheme:} AICTE Research Promotion Scheme (RPS). While undergraduate students are not eligible to serve as Principal Investigator (PI) under AICTE guidelines (which require a regular faculty member at an approved PG institution), this project provides a working preliminary prototype for a faculty mentor's future grant application in Educational Technology.

\section{Conclusion and Future Work}
This study presented the Cross-Venue Formatting Standardization System, an explainable, rule-based meta-template framework that bridges structural discrepancies across IEEE, ACM, and Springer conference templates. By formulating an 11-category canonical schema, engineering a deterministic mapping engine, and constructing an automated compliance validator, the system achieved a 95.56\% mapping accuracy and 100\% compliance detection on an independent 45-sample benchmark. The tool provides a transparent pedagogical scaffold for novice authors, reducing formatting overhead. Future work will explore hybridizing the rule engine with lightweight local sentence embeddings (such as SciBERT or MiniLM) to resolve unseen synonyms while maintaining complete explainability.

\section*{Acknowledgment}
The authors thank the Department of Computer Science for technical facilities and computational support during the execution of this Research Methodology project. This study was conducted as an undergraduate research project in academic paper formatting standardization.

\begin{thebibliography}{00}
\bibitem{jiang2019time} Y. Jiang, C. L. Gao, L. Y. Che, and X. X. Zhang, ``Time spent on article formatting: A survey among researchers,'' \textit{PLOS ONE}, vol. 14, no. 9, p. e0223976, Sep. 2019. DOI: 10.1371/journal.pone.0223976.
\bibitem{ahmed2020systematic} I. Ahmed and M. T. Afzal, ``A Systematic Approach to Map the Research Articles' Sections to IMRAD,'' \textit{IEEE Access}, vol. 8, pp. 129359--129371, Jul. 2020. DOI: 10.1109/ACCESS.2020.3009021.
\bibitem{ieee2024author} IEEE, ``IEEE Conference Template and Author Guidelines,'' IEEE Author Center, Piscataway, NJ, USA, 2024. [Online]. Available: \url{https://www.ieee.org/conferences/publishing/templates.html}
\bibitem{acm2024template} Association for Computing Machinery, ``ACM Primary Article Template and Publication Workflow (sigconf),'' ACM Digital Library, New York, NY, USA, 2024. [Online]. Available: \url{https://www.acm.org/publications/proceedings-template}
\bibitem{springer2024lncs} Springer Nature, ``Information for Authors of Computer Science Proceedings (LNCS),'' Springer-Verlag, Heidelberg, Germany, 2024. [Online]. Available: \url{https://www.springer.com/gp/computer-science/lncs/conference-proceedings-guidelines}
\bibitem{duransilva2024sasc} N. Duran-Silva, J. Moreno-Schneider, C. Parra-Rojas, and G. Rehm, ``Scientific Article Section Classification (SASC) Dataset,'' in \textit{Proc. 2024 Joint Int. Conf. Comput. Linguistics, Lang. Resour. Eval. (LREC-COLING 2024)}, Torino, Italy, May 2024, pp. 1656--1665. DOI: 10.63317/3rdo9r4az3iy.
\bibitem{lo2020s2orc} K. Lo, L. L. Wang, M. Neumann, R. Kinney, and D. S. Weld, ``S2ORC: The Semantic Scholar Open Research Corpus,'' in \textit{Proc. 58th Annu. Meeting Assoc. Comput. Linguistics (ACL 2020)}, Jul. 2020, pp. 4969--4983. DOI: 10.18653/v1/2020.acl-main.447.
\bibitem{teufel1999annotation} S. Teufel, J. Carletta, and M. Moens, ``An annotation scheme for discourse-level argumentation in research articles,'' in \textit{Proc. 9th Conf. Eur. Chapter Assoc. Comput. Linguistics (EACL 1999)}, Bergen, Norway, Jun. 1999, pp. 110--117. DOI: 10.3115/977035.977050.
\bibitem{sollaci2004introduction} L. B. Sollaci and M. G. Pereira, ``The introduction, methods, results, and discussion (IMRAD) structure: a fifty-year survey,'' \textit{J. Med. Libr. Assoc. (JMLA)}, vol. 92, no. 3, pp. 364--371, Jul. 2004.
\bibitem{councill2008parscit} I. G. Councill, C. L. Giles, and M.-Y. Kan, ``ParsCit: An open-source CRF reference string and logical document structure parsing package,'' in \textit{Proc. 6th Int. Conf. Lang. Resour. Eval. (LREC 2008)}, Marrakech, Morocco, May 2008, pp. 661--667.
\bibitem{lopez2009grobid} P. Lopez, ``GROBID: Combining Automatic Bibliographic Data Recognition and Term Extraction for Architecture and Engineering Research,'' in \textit{Research and Advanced Technology for Digital Libraries (ECDL 2009)}, LNCS, vol. 5743, Springer, Berlin, Heidelberg, Sep. 2009, pp. 473--474. DOI: 10.1007/978-3-642-04346-8\_62.
\bibitem{tkaczyk2015cermine} D. Tkaczyk, P. Szostek, M. Fedoryszak, P. J. Dendek, and \L. Bolikowski, ``CERMINE --- automatic extraction of structured metadata and references from scientific literature,'' \textit{Int. J. Doc. Anal. Recognit. (IJDAR)}, vol. 18, no. 4, pp. 317--335, Dec. 2015. DOI: 10.1007/s10032-015-0249-8.
\bibitem{jats2021} National Information Standards Organization, ``ANSI/NISO Z39.96-2021: Journal Article Tag Suite (JATS), version 1.3,'' NISO, Baltimore, MD, USA, Tech. Rep., 2021.
\bibitem{meuschke2023benchmark} N. Meuschke, V. Schubotz, F. Hamborg, T. Skopal, and B. Gipp, ``A Benchmark of PDF Information Extraction Tools Using a Multi-task and Multi-domain Evaluation Framework for Academic Documents,'' \textit{arXiv preprint arXiv:2303.09957}, Mar. 2023.
\bibitem{kashyap2023scientific} A. R. Kashyap, Y. Yang, and M.-Y. Kan, ``Scientific document processing: challenges for modern learning methods,'' \textit{Int. J. Digit. Libr.}, vol. 24, pp. 219--246, Aug. 2023. DOI: 10.1007/s00799-023-00352-7.
\bibitem{macfarlane2024pandoc} J. MacFarlane, ``Pandoc: A Universal Document Converter,'' Software manual, 2024. [Online]. Available: \url{https://pandoc.org}
\end{thebibliography}

\end{document}
